%% ====================================================================
\section{The LSS-Pulsar Adapter}
\label{sec:lss-pulsar}
%% ====================================================================

Pulsar's lifecycle (\S\ref{sec:lifecycle}) is concrete; the
\emph{orchestration} of generations, rollback, snapshots, and
dealer/coordinator role separation is provided by an external
universal-MPC framework. We use \textbf{LSS} \cite{luxlssmpc} ---
Lux's Linear Shamir Secret Sharing dynamic-resharing layer
\cite{lp077}, originally specified for ECDSA-CMP and FROST and
explicitly designed as a protocol-agnostic lifecycle extension.

The contract:
\begin{quote}\itshape
LSS owns generations, snapshots, rollback, and dealer/coordinator
roles. Pulsar owns lattice math: $R_q$ shares, lattice Pedersen
commits, $\Lambda$/Seeds/MAC regeneration, and the activation
certificate's lattice signature bytes.
\end{quote}

The wiring lives at \texttt{threshold/protocols/lss/lss\_pulsar.go} ---
the analogue of \texttt{lss\_cmp.go} (ECDSA) and \texttt{lss\_lens.go}
(curve, see LP-103 \cite{lp103}). The adapter exposes:

\begin{verbatim}
func DynamicResharePulsar(
    oldConfigs map[party.ID]*PulsarConfig,
    newPartyIDs []party.ID,
    newThreshold int,
    pool *pool.Pool,
    rand io.Reader,
) (map[party.ID]*PulsarConfig, error)
\end{verbatim}

It (1)~validates that all old configs agree on KeyEraID, Generation, and
GroupKey pointer; (2)~reconstructs an in-process \texttt{KeyEra} from
the union of old shares; (3)~delegates all lattice math to the Pulsar
kernel; (4)~bumps Generation by one, preserves KeyEraID, sets
RollbackFrom $= 0$ on forward transitions; (5)~persists the resulting
state through \texttt{PulsarSnapshotManager} for potential rollback.

\subsection{Bootstrap Dealer vs Signature Coordinator}

The LSS paper distinguishes two operational roles, both inherited by
Pulsar:

\begin{tabular}{@{}p{0.20\textwidth}p{0.36\textwidth}p{0.38\textwidth}@{}}
  \toprule
  \textbf{Role} & \textbf{LSS responsibility} & \textbf{Pulsar / Quasar mapping} \\
  \midrule
  Bootstrap Dealer & Network membership manager. Orchestrates resharing. Never holds unencrypted secrets. & Foundation MPC node at chain genesis (Bootstrap). P-Chain governance authority for Reanchor. Holds no shares after ceremony close. \\
  Signature Coordinator & Operational workhorse. Collects partial signatures. Combines certificates. & Quasar block proposer / per-epoch consensus runtime. Coordinates Sign1/Sign2/Combine for each Pulsar pulse over a bundle. \\
  \bottomrule
\end{tabular}

Crucially, the two roles are \emph{distinct authorities}: a compromised
Signature Coordinator cannot reshare or change the validator set
(it has no membership-management authority), and a compromised
Bootstrap Dealer at non-Bootstrap times cannot sign (it holds no
shares). The role split is the operational basis for Pulsar's
no-slashing-dependency failure model
(\S\ref{sec:lifecycle:rollback}).

\subsection{Adapter contract guarantees}

\textbf{Preserved across the call.} KeyEraID, GroupKey pointer, master
secret $\mathbf{s}$ (held only as shares; never reconstructed in any
process), and the persistent $(\mathbf{A}, \tilde{\mathbf{b}})$ public
key.

\textbf{Changed across the call.} Generation (incremented by 1),
participant set, threshold, share values, $\Lambda_i$ Lagrange
coefficients, pairwise PRF Seeds, MAC keys, transcript hash.

\textbf{Failure behavior.} Returns an error. The caller (LSS Rollback
Manager or Quasar consensus) decides whether to retry, rollback, or
reanchor. The adapter does not silently fall back.

\subsection{What the adapter does \emph{not} inherit}

A clean adapter must reject ECDSA-specific assumptions baked into
LSS-CMP and the curve-side LSS-FROST adapter (LP-103 \cite{lp103}):

\begin{itemize}[leftmargin=2em,nosep]
  \item ECDSA auxiliary secrets $w$, $q$ for nonce blinding (LSS
        Section 4 protocols I, II). Ringtail-style signing has Gaussian
        blinding built into Sign1/Sign2; no $w, q$ needed.
  \item Curve-point Pedersen commits $g^s \cdot h^r$. Pulsar uses
        lattice commits $\mathbf{A}\cdot\mathrm{NTT}(s) + \mathbf{B}
        \cdot\mathrm{NTT}(r)$ over $R_q$.
  \item ECDSA scalar serialization. Pulsar uses Lattigo \texttt{v7}
        polynomial serialization.
  \item Curve-specific Lagrange interpolation. Pulsar uses
        \texttt{primitives.ComputeLagrangeCoefficients} over $\mathbb{Z}_q$.
\end{itemize}

\subsection{Acceptance tests}

Ten acceptance tests gate the adapter (collapsed into six
\texttt{Test*} functions):

\begin{enumerate}[leftmargin=2em,nosep]
  \item GroupKey pointer preserved across resharing.
  \item KeyEraID preserved across resharing.
  \item Generation increments by exactly one on forward transition.
  \item RollbackFrom $= 0$ on forward transition.
  \item $t_{\text{old}} \neq t_{\text{new}}$ works.
  \item Old set $\neq$ new set works (disjoint rotation supported).
  \item Regenerated shares produce a valid Pulsar threshold signature.
  \item Signature verifies under \emph{unchanged} GroupKey, proving
        old shares are not needed after activation.
  \item Pairwise PRF / MAC material is regenerated for exactly the new
        party set, with bytes that differ from the old committee's.
  \item Rollback restores prior generation snapshot; restored state
        carries Generation $= \text{current} + 1$ and RollbackFrom
        $= \text{current}$, preserving monotonicity while marking
        lineage.
\end{enumerate}

All ten pass at the time of writing
(\texttt{threshold/protocols/lss/lss\_pulsar\_test.go}).
